%%%%%%%%%%%%%%%%%%%%%%%%%%%%%%%%%%%%%%%%%%%%%%%%%%%%%
%
%           EE474 Hardware and Basic I/O
%
   % combined into notes pack	%<stream, c>

\documentclass{article}	%<nh>

\usepackage{graphicx}	%<shn>
\usepackage{hyperref}	%<shn>
\usepackage{listings}	%<shn>

	%<*hn>
\oddsidemargin 0in
\evensidemargin 0in
\textwidth 6.5in
\textheight 8.7in
\topmargin 0.2in


\linespread{1.0}     % 1.0 = single space   1.6 = double space


	%<*>
\title{EE 472:  Serial Communication}

\author{Blake Hannaford\\
        Department of Electrical \& Computer Engineering \\
        The University of Washington
}

\date{\today}

\begin{document}	%<shn>


 

\maketitle
\subsection*{INSTRUCTOR NOTES}	%<n>

	%<*>
%%%%** Section 1 
\section{Serial Communication}

Send bits from point A to point B one at a time.

Info Sources:
{\tt http://www.seetron.com/ser\_an1.htm} \\
{\tt http://www.taltech.com/TALtech\_web/resources/intro-sc.html }\\
(on-line 12-Nov-02)

\section*{RS232C}

An international standard for serial communication (~1960s).

Very old but still influential in ~2020 systems.

%%%%** Section 1.1
\subsection{Electrical Signaling}

\begin{tabular}{|l|c|c|}\hline
+2.5V --- +15V  &  logic 0 &  ``space'' \\ \hline
-2.5V --- -15V  &  logic 1 &  ``mark''  \\ \hline
\end{tabular}

High voltage {\it margin} gives some immunity to noise.

Also a hassle since many systems don't otherwise
need +15V, -15V supplies.

Requires special driver chips for +- voltages.

Today some systems use

\begin{itemize}
        \item +5V logic 0
        \item  0V logic 1
\end{itemize}

to save power supply costs.


%%%%** Section 1.2
\subsection{Serial Signaling}

{\it Idle} line in ``mark'' state (logic 1, -15V)

\begin{tabular}{|c|c|l|}
\hline
Pos.    & Value         & Function      \\ \hline
0       &  space        & Start Bit     \\ \hline
1       &  X            & Data 0 (lsb)  \\ \hline
2       &  X            & Data 1        \\ \hline
...     &  X            & ...           \\ \hline
8       &  X            & Data 7 (msb)  \\ \hline
9       &  X            & Parity*       \\ \hline
10      &  mark         & Stop Bit 1    \\ \hline
11      &  mark         & Stop Bit 2*   \\ \hline
\end{tabular}

\noindent
{\bf * Note:} Parity and Stop Bit 2 are optional,
can have 1.5 Stop
Bits.


\includegraphics[width=\linewidth]{Serial_figs/rs232_serialbits.png}

\subsection{Parity}

\begin{verbatim}
P = D0 ^ D1 ^ D2 ... ^ D7
\end{verbatim}
(XOR of all data bits).

Interpretation of XOR(bits) is
\begin{itemize}
   \item "1" if number of ones is ODD 
   \item "0" if number of ones is EVEN
\end{itemize}

Example: 

\begin{tabular}{|c|c|c|c|}\hline
D0 & D1 & D2 & XOR(D0,D1,D2) \\ \hline
0  &  0 & 0  &  0 \\\hline
0  &  0 & 1  &  1 \\\hline
0  &  1 & 0  &  1 \\\hline
0  &  1 & 1  &  0 \\\hline
1  &  0 & 0  &  1 \\\hline
1  &  0 & 1  &  1 \\\hline
1  &  1 & 0  &  0 \\\hline
1  &  1 & 1  &  1 \\\hline
\end{tabular}

Check if received correct parity bit. 


%%%%** Section 1.3
\subsection{Speed and Timing}

{\bf Baud} is short for Baudot, 19th century
French telegrapher who invented first fixed length
alphanumeric code.

{\bf Baud} means number of signal changes on the
line per second.  Commonly equal to bits per second.

{\bf Example: }  Arduino serial port default: 9600 Baud. (it's actually USB!)


%%%%** Section 1.4
\subsection{Synchronization}
There is no common clock between sender and
receiver. How does receiver know when to test
levels?

\noindent
{\tt A1}: Synchronus mode --- rarely used.  Continuous
stream of characters.

\noindent
{\tt A2}: Asynchronus mode --- start bit is known
to be logic 0.  Falling (rising?) edge triggers
a local clock in the receiver.


%%%%** Section 1.5
\subsection{Multiplexing}
Sending information in both directions.  (Stringing 2 telegraph wires was very expensive!)

{\bf Terms:}

{\it Simplex}  One sender and one receiver.

{\it Half Duplex} One side can send at a time.

{\it Full Duplex} Both sides can send at same time.

Typical serial communication today is {\it Full Duplex}.
%%%%** Section 1.6
\subsection{RS232 Wiring}

Minimum 3 wires:
\begin{itemize}
        \item TX  (transmit)
        \item RX  (receive)
        \item gnd  (V=0 reference)
\end{itemize}

Maximum 23 wires!

\includegraphics[height=2.75in]{Serial_figs/rs232_db25pinout.png}
{\tiny
Source(https://www.aggsoft.com/rs232-pinout-cable/}
{\tiny serial-cable-connections.htm)}

%%%%** Section 1.7
\subsection{Connections}
Must have:

\noindent
TX $\to$ RX \\
RX $\gets$ TX

Cables are easier to make if pins are
``straight through''. i.e. pin {\it n}
connects to pin {\tt n}.

%%%%** Section 1.8
\subsection{Connectors}
23 pin (ancient, un-needed)

9 pin (IBM PC standard)

Two connector wiring types

{\tt DTE:}  pin 2 = RX, pin 3 = TX ``Data Terminal Equipment"

{\tt DCE:}  pin 2 = TX, pin 3 = RX ``Data Communications Equipment"

Thus, {\tt DCE }can talk to {\tt DTE} with
a ``straight through'' cable.


{\tt DTE, DCE} are ancient
names from pre-PC days!!.

Typical 1990's:

PC == {\tt DTE } \\
modem, printer, etc, == {\tt DCE}


%%%%** Section 1.9
\subsection{Multi-Drop}

Can also have more than 2 systems interconnected
by a {\it serial bus}.

In that case all users transmit and listen to
same wire.

Must implement protocol to prevent two or more
users from talking simultaneously.

(like half-duplex)

%%%%** Section 1.10
\subsection{``Null Modem''}

What if you want to connect a {\tt DTE} to
another {\tt DTE} (such as PC-PC)?

Use a special cable which connects all lines
``straight through'' {\it except}

pin 2 $\to$ pin 3 {\it and} \\
pin 3 $\to$ pin 2


%%%%** Section 1.11
\subsection{Flow Control}

Sometimes a slow device cannot handle data as
fast as sender (even at same baud rate).

Receiver needs a way to stop and start sender.

{\bf Hardware solution:} Additional wires:\\
{\tt RTS} ``Request to Send".  {\tt DTE}
sets this to mark to allow data to be sent to
it by {\tt DCE}. \\
{\tt CTS} ``Clear to Send". {\tt DCE} sets this
to mark to allow data to be sent to it by {\tt DTE}.

Hardware flow control can be enabled/disabled.

{\bf Software solution:} Special Control Characters:\\
{\tt XON} ({\tt ctl Q}) Start Transmitting\\
{\tt XOFF}({\tt ctl S}) Stop Transmitting



%%%%** Section 1.12
\subsection{Flow Control Note}

{\bf Note:}  If you use Hardware Flow control,
then ``null modem'' cable must also cross the
{\tt RTS } and {\tt DTS} wires (pins 7,8 on 9-pin
IBM standard).



%%%%** Section 2 
\section{Error Control}
%%%%** Section 2.1
\subsection{Causes}
Unfortunately, serial communication is subject
to errors.  Sources:

\begin{itemize}
        \item Electrical interference
        \item Long wire lengths
        \item Ground potential differences between
        communicating systems.
        \item Timing errors between sender and
        receiver clocks (should be within 1\%).
\end{itemize}


%%%%** Section 2.2
\subsection{Solutions}
\begin{description}
 \item[Character Parity] 9th data bit.
        Very weak and bandwidth hog.
 \item[Checksum]  Add up a group of bytes.
        Send sum at end of packet. Compare.
 \item[Parity] Send parity of some group of
        bits (rows or colums) in the packet.
 \item[CRC] Use theory of binary polynomials.
       Send a code so that combined message
       is divisible by some known polynomial.
\end{description}

%%%%** Section 2.3
\subsection{Packets}

Grouping bytes into packets gives structure
to the serial communication.   Packet structure
must be same between software both sides.

Typical Packet:


{\tt
\begin{tabular}{|c|c|c|c|c|}
\hline
Start &   Type & Length  & ... & Checksum  \\
\hline
\end{tabular}
}

\begin{description}
        \item[Start] A character which the software can recognize
        as the start of a packet.   Ideally should not appear anywhere
        else in packet.
        \item[Type] Kind of packet.
        \item[Length] value which tells how long the
        packet will be.
        \item[...] some number of data bytes (the payload).
        \item[Checksum] Checksum computed on previous packet.
\end{description}

%%%%** Section 2.4
\subsection{Packet Error Control}

What if a packet is received and the Checksum is wrong?

\noindent
Normal:

1. Sender sends the packet and {\it saves} the packet.\\
2. Receiver checks Checksum and {\it acknowledges} receipt of packet.\\
3. Sender discards packet and sends next packet.

\noindent
Error:

1. Sender sends and {\it saves} the packet.\\
2. Receiver detects error and sends {\it negative acknowledge}.\\
3. goto step 1.


%%%%** Section 2.5
\subsection{ACK/NAK}

Special short packets:

{\tt ACK} Acknowledge
{\tt NAK} Negative Acknowledge

{\tt
\begin{tabular}{|c|c|c|}
\hline
Start &   ACK/NAK &  Checksum  \\
\hline
\end{tabular}
}

%%%%** Section 2.6
\subsection{Protocol Issues}

What Happens If:
\begin{itemize}
        \item Transmission is so unreliable that Checksum fails
                every time?
        \item {\tt ACK/NAK} packets are subject to errors?
\end{itemize}
Do we have to send {\tt ACK/NAK} on receipt of an {\tt ACK/NAK}?

Answers to all these questions must be worked out in a
protocol spec.

\end{document}	%<shn>

